% !TeX program = xelatex
% 编译方式：使用 XeLaTeX 编译（ctex 宏包需要 XeLaTeX 或 LuaLaTeX）
\documentclass[11pt,a4paper]{ctexart}

\usepackage[margin=2.4cm]{geometry}
\usepackage{amsmath,amssymb}
\usepackage{booktabs}
\usepackage{array}
\usepackage{graphicx}
\usepackage{xcolor}
\usepackage{enumitem}
\usepackage{listings}
\usepackage{tikz}
\usetikzlibrary{arrows.meta,positioning}
\usepackage[hidelinks]{hyperref}
\usepackage{caption}

\captionsetup{font=small,labelfont=bf}
\setlist{itemsep=2pt,topsep=3pt,parsep=0pt}

\definecolor{codebg}{RGB}{246,248,250}
\lstset{
  basicstyle=\ttfamily\small,
  backgroundcolor=\color{codebg},
  frame=single,
  rulecolor=\color{lightgray},
  breaklines=true,
  showstringspaces=false,
  columns=fullflexible,
  xleftmargin=6pt,xrightmargin=6pt,
  aboveskip=8pt,belowskip=6pt,
}

\title{\bfseries 面向灾害新闻的信息抽取系统\\[2pt]\large （信息与知识获取 · 作业三 实验报告）}
\author{}
\date{2026 年 6 月}

\begin{document}
\maketitle

\begin{center}
\small
\begin{tabular}{>{\centering\arraybackslash}m{2.2cm}>{\centering\arraybackslash}m{2.8cm}>{\centering\arraybackslash}m{1.4cm}>{\raggedright\arraybackslash}m{5.6cm}}
\toprule
姓名 & 学号 & 角色 & 分工 \\
\midrule
曹忠昊 & 2023211964 & 组长 & 灾害语料爬虫与三重过滤、系统整合 \\
侯钧译 & 2023211728 & 组员 & 集成抽取算法（正则 / 触发词 / NER 加权投票）与知识图谱 \\
徐舸山 & 2023211038 & 组员 & EasyOCR 多媒体抽取、人工评价与实验报告 \\
\bottomrule
\end{tabular}
\end{center}

\vspace{0.4em}

\begin{abstract}
\noindent
本作业实现了一个面向真实灾害新闻的信息抽取系统。系统从政府与通讯社等权威来源采集了 118 篇真实灾害报道并本地存储，
对每篇报道抽取灾害类型、发生时间、发生地点、伤亡情况、经济损失、响应级别与救援信息这 7 个信息点，并把它们组装成一个结构化的灾害事件。
抽取方法以正则匹配为基础，在此之上做了集成：正则、触发词与基于 jieba 词性标注的命名实体识别三路并行，按字段加权投票，
消融实验显示集成相对纯正则把平均填充字段数提升了约 4.8\%。事件进一步聚合为知识图谱。
在文本之外，系统先用视觉大模型（Qwen-VL）对配图做质检、剔除 logo、二维码、UI 等无关图，再对保留的真实灾情图用 EasyOCR 识别文字并回灌进抽取，是图像与文本结合的多媒体信息抽取尝试。
抽取结果支持逐字段的人工评价。除从零跑 OCR 需要首次联网下载模型外，全流程在 CPU 上运行，且不含任何编造数据。
\end{abstract}

\section{任务背景与总体设计}

信息抽取要做的，是从非结构化的文本里把我们关心的特定信息点抽出来，整理成结构化的形式。本次作业的基本要求是：
选定一个特定领域，用爬虫采集不少于 100 篇文档并本地存储，对感兴趣的信息点做抽取并展示，信息点不少于 5 个且能组合成一个事件，
可以调用开源的中英文 NLP 基本模块，抽取算法至少要实现正则表达式匹配，并尽量对结果准确率做人工评价；
扩展要求鼓励尝试多媒体抽取与算法优化。

我们选的领域是灾害事件。一条灾害新闻天然适合做信息抽取：它通常会交代是什么灾、发生在何时何地、造成了多少伤亡和损失、
启动了几级应急响应、有哪些救援力量参与。我们把这些归纳成 7 个信息点，它们恰好能拼成一个完整的灾害事件，
也超过了「不少于 5 个」的要求。在抽取算法上，除了必需的正则匹配，我们还加了触发词和命名实体识别两路，做成加权投票的集成，
这是本系统主要的算法投入；NLP 基本模块用的是 jieba 的分词、词性标注与分句。系统的整体流程见图~\ref{fig:arch}。

\begin{figure}[ht]
\centering
\begin{tikzpicture}[
  node distance=8mm and 11mm,
  box/.style={draw,rounded corners=2pt,align=center,inner sep=4pt,minimum height=8mm,font=\small},
  >={Stealth[length=2mm]}
]
\node[box] (crawl) {灾害爬虫\\(三重过滤)};
\node[box,right=of crawl] (corpus) {灾害语料 + 配图};
\node[box,right=of corpus] (ocr) {VLM 质检\\+ EasyOCR};
\node[box,right=of ocr] (ext) {集成抽取\\正则+触发词+NER};
\node[box,below=10mm of ext] (event) {7 要素结构化事件};
\node[box,left=of event] (kg) {知识图谱};
\node[box,left=of kg] (eval) {逐字段人工评价};
\draw[->] (crawl) -- (corpus);
\draw[->] (corpus) -- (ocr);
\draw[->] (ocr) -- (ext);
\draw[->] (ext) -- (event);
\draw[->] (event) -- (kg);
\draw[->] (event) -- (eval);
\end{tikzpicture}
\caption{系统流程：爬虫采集灾害语料，配图经 VLM 质检与 OCR 后文字回灌，集成抽取产出 7 要素事件，再聚合为知识图谱并支持逐字段人工评价。}
\label{fig:arch}
\end{figure}

% 如需插入系统界面截图，取消下面注释并放入图片文件：
% \begin{figure}[ht]\centering
% \includegraphics[width=0.9\textwidth]{screenshot_extract.png}
% \caption{抽取结果页：7 个信息点逐字段给出取值、集成置信度与来源标注，并附人工评价按钮。}
% \end{figure}

\section{灾害语料采集}

灾害领域的语料采集不能只追求数量，更要保证「确实是灾害事件」。我们在通用异步爬虫的基础上做了主题聚焦和质量过滤，
最终入库 118 篇真实灾害报道，去重后覆盖 19 个域名、14 个来源，全部是政府与通讯社等权威源。
来源较多的有四川省政府（22）、应急管理部（22）、水利部（12）、人民网（12）、新华网（10）、中国地震局（8）、
中国新闻网（7）、澎湃新闻（7）；按灾害类型分，地震 39、洪水 33、暴雨 25、台风 11、火灾 5、山体滑坡 3、暴雪 2，共 7 类。
每篇文档都保留原始 URL、来源与日期，可逐条溯源；图片路径统一规范化成 \texttt{data/images/} 下的相对路径，
不含机器绝对路径，整个项目可以直接迁移；代码里也删除了一切数据生成器，没有伪造文档。

值得一提的是过滤这一环，因为它直接决定了语料的「事件纯度」。最初版本从地震局等站点的链接里带回了大量机构话术页面——
党组读书班、警示教育、巡视通报、座谈调研、监测站简介、信息公开年报、网站地图等，这些页面字面上含灾害关键词，
却根本不是灾害事件。为此我们在原有的灾害特征判断和真实性校验之外，加了一道事件性过滤：要求标题里出现具体的灾害事件或响应信号
（震级、启动响应、遇难、受灾、登陆、强降雨、防汛、救灾资金等），同时把党务、会议、解读、年报、导航、人名这类机构话术标题挡在外面。
配合来源白名单和单源配额，最终留下的才是真实灾害事件而非部门工作动态。爬虫本身仍是异步并发、礼貌抓取、SimHash 去重，
配图按 PNG/JPG 文件头校验后落盘，为后面的 OCR 提供素材。

\section{信息抽取方法}

\subsection{七个信息点}

系统抽取的 7 个信息点及其主要抽取手段见表~\ref{tab:fields}。需要澄清的是，「信息点不少于 5 个」指的是系统定义并抽取的信息点种类，
也就是事件模板的槽位数，这里是 7 个；至于某一篇报道最终填满几个，取决于这条新闻本身写了多少，这一点会在实验结果里如实讨论。

\begin{table}[ht]
\centering
\caption{七个信息点与主要抽取手段}
\label{tab:fields}
\begin{tabular}{lll}
\toprule
信息点 & 含义 & 主要抽取手段 \\
\midrule
灾害类型 & 地震 / 台风 / 洪水 / 暴雨 / 火灾 …… & 正则 + 触发词 \\
发生时间 & 日期 / 时段 & 正则 + NER（时间） \\
发生地点 & 省市县 / 震中 & NER（地名）+ 正则 \\
伤亡情况 & 死亡 / 受伤 / 失踪人数 & 正则 \\
经济损失 & 直接经济损失金额 & 正则 \\
响应级别 & Ⅰ--Ⅳ 级应急响应 & 正则 \\
救援信息 & 转移 / 搜救 / 救灾力量 & 触发词 + NER（机构） \\
\bottomrule
\end{tabular}
\end{table}

\subsection{集成抽取与加权投票}

单靠正则有两个毛病：表述一灵活就漏抽（召回不足），有时又会抽到不该抽的（精度受限）。所以我们把三个互补的抽取器组合起来，
按字段加权投票。正则引擎模式精确，主攻响应级别、伤亡数字、经济损失金额这类格式很强的字段，权重基线设得较高（0.6）；
触发词引擎基于触发词加上下文窗口，补正则覆盖不到的自然语言表述；NER 引擎用 jieba 的词性标注（\texttt{jieba.posseg}）
识别人名、地名、机构和时间，再用领域规则补全完整的行政区划与救援机构，主要补地点、机构和时间。

这里有一个实现上的取舍值得说明。我们原本打算用百度的 LAC 做 NER，但实际部署环境里并没有装 LAC，
如果硬依赖它，NER 这一路就会静默退化成正则兜底，报告里写的「用了 NER」就名不副实了。权衡之后，我们改用随处可得的
\texttt{jieba.posseg} 作为默认的开源 NER，把 LAC 降级为「装了才用」的可选增强。这样无论在什么环境下，
NER 这一路都确实在调用一个开源 NLP 模块，而不是挂个名头。jieba 的词性标注对机构名和细粒度地名召回偏弱，
这部分正好由领域规则补上，二者互补。

投票的具体流程是：对每个字段先收集各抽取器给出的候选，按归一化和子串包含关系把同义候选合并成组，再取累计权重最高的那一组。
若记字段 $f$ 上产生了候选的抽取器集合为 $M_f$、胜出候选所在的组为 $G^\ast$，则该字段的集成置信度定义为
\begin{equation}
\text{conf}(f)=\frac{\sum_{m\in G^\ast} w_{f,m}}{\sum_{m\in M_f} w_{f,m}},
\end{equation}
即胜出组的权重占所有发声抽取器权重的比例。每个字段最终都带着这个置信度和来源标注（如 \texttt{regex+ner}）展示在网页上。

举个发生地点的例子（该字段三个抽取器的权重为正则 0.45、触发词 0.30、NER 0.40）。若正则与 NER 都抽到「甘肃临夏州积石山县」，
两者归并为同一组、权重合计 0.85，没有竞争候选，置信度即 $0.85/0.85=1.0$；但若正则抽到「甘肃临夏州积石山县」、NER 却抽到「青海」，
两组权重为 0.45 与 0.40，前者胜出，置信度为 $0.45/0.85\approx0.53$——这个偏低的值恰好提示该字段存在分歧、值得人工复核。
页面再据此给字段着色（绿 $\ge 0.8$、橙 $\ge 0.5$、灰更低），相当于替人工评价排好了优先级：绿色的可以快速放过，橙、灰的重点核对。
对 NER 常见的边界噪声（比如把「当日」「日从」之类前缀粘进实体），我们用规则做了清洗；
发生地点还专门做了行政区划规范化——以省份名锚定抽取最长的行政链，剥掉「在 / 家住 / 造成」这类非地名引导词，
合并重复的行政后缀，并剔除「重点地区 / 南方地区」这种泛化表述，宁可少给也不给脏数据。

\subsection{消融实验}

为了量化集成相对纯正则的收益，我们在同一份语料上对比了两者的平均填充字段数：纯正则为 3.70，集成为 3.88，提升约 4.8\%。
增益主要集中在发生地点、救援信息、伤亡情况这几个「软」字段上——这些字段表述灵活，靠 NER 补地名机构、靠触发词补自然语言说法；
尤其是发生地点，在正则被规范化严格收紧之后，干净候选更多来自触发词和 NER。而响应级别、灾害类型这类强格式字段，
正则本就能高精度命中，集成在这里的作用更多是加权去噪、给出置信度。

\section{事件组装与知识图谱}

抽取出来的字段会被组装成一个结构化的灾害事件，附带严重程度定级、信息完整度、缺失字段列表和一句摘要。
进一步地，我们把这些事件聚合成知识图谱：节点包括事件、地点、灾害类型、机构、时间等，边把事件和它的各个要素连起来。
当前图谱共 386 个节点、641 条关系，其中事件节点 105 个（即 107 篇有效事件中有 105 篇要素足够充分、成功建立了事件节点）。
前端用 ECharts 的力导向图展示，可以按类型或地区做态势上的观察。

\section{多媒体信息抽取：VLM 图片质检与 OCR}

灾害新闻页抓来的「配图」里混着大量与灾情无关的噪声——站点 logo、二维码、UI 图标、广告横幅、纯文字海报等。
如果不加甄别直接 OCR，既会让配图展示一团乱，也会让「图片文字」里塞满从 logo、二维码识别出的乱码。所以配图我们分两步处理。

\subsection{VLM 图片质检}

把每张配图降到 512px，送进视觉大模型（通义千问 \texttt{qwen3-vl-plus}，经 DashScope 的 OpenAI 兼容接口），
让它判断这是不是一张「有意义的灾害新闻内容图」，只输出 \texttt{\{keep, type, reason\}}。判定结果按文件名缓存到
\texttt{data/vectors/vlm\_verdicts.json}，随交付一起分发——即便批改环境没有密钥，也能复现「已清洗」的结果；
判定失败则 fail-open 默认保留，密钥只从环境变量读取、不写入任何文件。全库 91 张配图送检，剔除了 71 张：
网站 logo 25、二维码 10、装饰图 7、UI 图标 6、与灾害无关的文字海报 5 等，最终保留 20 张真实灾情内容图（分布在 12 篇文档里）。
这套做法与作业二的图片质检同源，只是把提示词换成了灾害题材。

一个有点出乎意料、但值得如实写下的发现是：权威灾害通报页里的「配图」大约八成是站点 logo、二维码、UI 这类页面装饰，
而不是灾情照片。VLM 质检把它们和真正的内容图干净地分开了。

\subsection{OCR 识别与回灌}

对\textbf{通过质检}的内容图，我们用 EasyOCR（\texttt{ch\_sim+en}，纯 CPU，置信度阈值 0.3）识别文字，
再把 OCR 文本以「\texttt{【图片文字】}」的形式拼接到正文之后，一起参与集成抽取。质检之后有 8 篇文档获得了
真正来自灾情图的 OCR 文本。识别结果缓存到 \texttt{ocr\_cache.json}，演示时秒开；OCR 页面还提供「实时重新识别」按钮，
可以现场验证识别确实是在线跑出来的，而不是预置的。

这里要补一个诚实的对照：把无关图的 OCR 文本从抽取输入里去掉后，7 要素的抽取结果\textbf{完全没变}
（有效事件、平均填充、各字段填充率都和去掉前一致）。这说明灾情要素本来就出自正文，而先前「82 篇参与 OCR」里
绝大多数其实是在识别 logo、二维码这些噪声。去掉之后「参与 OCR」从 82 降到 8，是更诚实的口径，
知识图谱也相应从 388/643 微调到 386/641。

实现上还踩过一个不太好查的坑：EasyOCR 底层用 \texttt{cv2.imread} 读图，而它在 Windows 上读不了中文路径，
会静默返回空、随后报 \texttt{NoneType} 没有 \texttt{shape}。因为我们的项目路径里带中文，OCR 一度整体失效。
解决办法是改用 \texttt{PIL.Image} 读图、转成 numpy 数组再喂给 EasyOCR，绕开 OpenCV 对中文路径的限制。

\section{实验结果}

\subsection{抽取覆盖率}

118 篇文档上各字段的填充率见表~\ref{tab:cover}。平均填充 3.88 个字段，有 107 篇（91\%）凑齐了关键要素、算作有效事件。

\begin{table}[ht]
\centering
\caption{各信息点在 118 篇上的填充率}
\label{tab:cover}
\begin{tabular}{lcl}
\toprule
信息点 & 填充率 & 说明 \\
\midrule
灾害类型 & 118/118（100\%） & 入库即要求命中灾害类型 \\
发生时间 & 103/118（87\%） & 正则多模式 + NER（时间） \\
发生地点 & 85/118（72\%） & NER 与正则协同，经规范化只保留具体行政区划 \\
救援信息 & 80/118（68\%） & 触发词 + 机构；权威源重点报道处置与救援力量 \\
响应级别 & 49/118（42\%） & 启动应急响应的事件含；支持中文「四级」$\to$ Ⅳ级归一化 \\
伤亡情况 & 14/118（12\%） & 含具体伤亡数字的事件 \\
经济损失 & 9/118（8\%） & 披露直接经济损失金额的事件 \\
\bottomrule
\end{tabular}
\end{table}

这张表里有两点值得展开。其一，伤亡和经济损失的填充率明显偏低，这不是抽取的缺陷，而是真实信息源的特点：
权威媒体的即时通报往往先讲清楚类型、时间、地点和响应级别，具体的伤亡和损失数字通常要等后续的跟踪报道才披露。
我们认为如实呈现这一现象比硬凑数字更诚实，也更符合「不编造数据」的要求。其二，发生地点的 72\% 是收紧后的结果。
在做行政区划规范化之前，地点字段的填充率看着有 94\%，但里面混着「家住某县县」「部署重点地区」「南方地区」这类噪声；
规范化之后收敛到 72\%，但留下的全是干净、具体的行政区划。我们宁愿用一点召回换取精度，因为在人工评价里，
脏地点会被直接判错，得不偿失。

\subsection{人工评价}

抽取结果页的每个字段旁都有「正确 / 部分 / 错误」三个按钮，评分者点击后分别按 1.0 / 0.5 / 0.0 计分，
系统实时返回这次的准确率和逐字段统计并写入评价库，可在汇总页按字段聚合查看。交付的版本里评价库是干净的初始状态，
分数由评分者现场打出来，是真实记录而非预置。

\section{对环境与社会可持续发展的影响}

环境方面，爬虫的异步并发与礼貌抓取减少了无效请求，SimHash 去重避免了重复存储，抽取结果、OCR 与图谱都是一次构建、
持久化复用，加上用的是 EasyOCR、jieba 这类 CPU 友好的模型，整体算力和能耗都比较低。社会方面，
把散落在新闻里的灾情要素自动结构化、再关联成事件图谱，可以为应急响应和防灾减灾提供线索，服务公共安全；
OCR 把图片里的灾情信息也补了进来，提升了应急信息的完整度。系统内设 \texttt{/sustainability} 页面对此做了说明。

\section{局限与展望}

伤亡和经济损失的填充率受信息源披露习惯的限制，后续可以做跨文档的事件归并，把同一场灾害的即时通报和后续跟踪报道关联起来补全数字。
NER 目前是 jieba 词性标注加领域规则，灾害领域的专名（细粒度地名、救援机构）还可以靠领域词典、微调，或在有条件时接入 LAC、HanLP 来增强。
知识图谱也可以再往前走一步，做事件去重和时序链，支撑同一灾害多篇报道的聚合。

\section{运行与复现}

系统已附带抽取结果、知识图谱与 OCR 缓存，开箱即可启动：

\begin{lstlisting}[language=bash]
pip install -r requirements.txt
python run.py web        # 启动 Web，默认 http://127.0.0.1:8090
# 从零构建（可选）：
python run.py crawl --max 400   # 异步爬取真实灾害新闻
python run.py ocr               # EasyOCR 识别配图文字并回灌
python run.py extract           # 全量集成抽取 + 构建知识图谱
\end{lstlisting}

\noindent
启动后可访问抽取、知识图谱、OCR 抽取、抽取评价与可持续等页面。NER 默认使用 jieba，无需联网下载；
仅从零跑 OCR 或「实时重新识别」时，首次会联网下载 EasyOCR 模型，之后离线复用。

\end{document}
